\subsection{Checklist: Declaration of LLM Usage}
\label{app:checklist_llm_usage}

\paragraph{Question.}
Does the paper describe the usage of LLMs if they are an important, original,
or non-standard component of the core methods in this research, while
distinguishing any writing, editing, or formatting assistance from the
scientific methodology?

\paragraph{Answer.}
Yes.

\paragraph{Justification.}
The current manuscript describes the proposed method as a sparse Koopman
autoencoder: an encoder, linear decoder, latent Koopman transition, sparse
training objective, and post-hoc support-family diagnostics. The paper does not
describe an LLM as part of the learned predictor, training loss, inference
procedure, rollout evaluation, support-family construction, or statistical
testing protocol.

The manuscript does disclose LLM involvement in benchmark construction. In the
multibasin benchmark appendix, the \(15\) two-dimensional multibasin benchmark
systems are described as having been generated procedurally with Claude Opus
4.5. Current source inspection also finds Claude-specific benchmark catalog
adapters, manifests, plotting utilities, validation utilities, and queueing
scripts. These artifacts support benchmark definition and experiment
infrastructure; they do not make an LLM a runtime component of the SKAE models
or of the reported forecasting and support-alignment evaluations.

\paragraph{Distinction between core methodology and assistance.}
\begin{itemize}
  \item \textbf{Core methodology: No LLM component.}
  The core method is the sparse Koopman autoencoder and its associated
  support-based evaluation. Model training and evaluation use numerical
  trajectory data, PyTorch model components, deterministic benchmark code, and
  post-hoc metrics. No LLM is used to encode states, choose supports, advance
  rollouts, assign support families, compute losses, select basin labels, or
  produce the reported experimental measurements.

  \item \textbf{Benchmark construction: LLM-assisted.}
  Claude Opus 4.5 was used to procedurally generate the multibasin benchmark
  systems reported in the benchmark inventory. These systems are subsequently
  represented as explicit dynamical-system code and benchmark metadata. Basin
  labels and basin counts are used only as benchmark metadata for evaluation and
  diagnostic slicing, except where the manuscript explicitly notes their use to
  set fixed architecture group counts for structured-transition ablations.

  \item \textbf{Writing and editing assistance: LLM-assisted.}
  LLM assistance was used to draft and edit this checklist declaration. Such
  assistance is language-level documentation support and is separate from the
  scientific method, model implementation, experiment execution, and numerical
  results.
\end{itemize}

\paragraph{Caveats.}
\begin{itemize}
  \item This declaration is based on the current manuscript and current code and
  documentation search, with the manuscript treated as the source of truth.

  \item The search confirms named Claude-related benchmark and infrastructure
  artifacts, but it is not a broader provenance audit of unpublished drafting
  records or external notebooks.

  \item The LLM-assisted benchmark-generation disclosure should be retained in
  any final submission version that uses these \(15\) multibasin systems. If
  additional LLM assistance is used later for manuscript text, code generation,
  experiment design, or artifact preparation, this declaration should be updated
  before submission.
\end{itemize}
